\specialchapter{术语索引}

参考编号依次表示章节、段落和编号（或有时表示习题）

实或复李群的伴随群: III.6.4. 李代数中元素的线性映射: I.1.2. 李代数的表示: I.3.1. 李群的表示: III.3.12.

阿多定理: I.7.3.

代数（不一定结合）: I.1.1. 由代数扩张标量得到的代数: I.1.1. 李代数的包络代数: I.2.1. 李代数: I.1.2. 相反代数: I.1.1. 乘积代数: I.1.1. 商代数: I.1.1. 李 p-代数的受限包络代数: I.2, 习题 6. 模的对称代数: I.2.5.

几乎单李群: III.9.8.

n 次交错式: II.2.6.

李代数包络代数的主反自同构: I.2.4.

与 g-模（一个表示）相关的（双线性型）: I.3.6. 与运算律相关的（无穷小运算律）: III.3.7.

自由结合代数: II.3.

李代数的特殊自同构: I.6.8.

基本交换子: II.5.4. 李代数的族: II.2.3.

比伯巴赫定理: III.4, 习题 13.

双代数: II.1.2.

双不变截面: III.3.13.

二项式多项式: II.5, 习题 4.

括号: I.1.2.

左典范微分形式: III.3.13, III.3.18. 卡当判据: I.5.4. Casimir 元: I.3.7. 中心扩张: I.1.7. 群的滤过: II.4.4. 中心化子: I.1.6, III.9.3. 李代数的中心: I.1.6. 特征理想: I.1.4. 幂零类: II.2.7. 简单表示的类: I.3.1. 取值于 g-模的余边界、余链、余循环: I.3, 习题 12. 余代数: II.1.1. 交换李代数: I.1.3. 相容的群结构和流形结构: III.1.1. 完全不变双线性型: I.3.6. 可约表示: I.3.1. 复李群: III.1.1, III.8.1. 实李群的复化: III.6.10. g-模的同型分量: I.3.1. 半单李代数的简单分量: I.6.2. 复李群的共轭李群: III.1.1. 李代数包络代数元素的常数项: I.2.1. Magnus 代数元素的项: II.5.2. 代数相对于基的结构常数: I.1.1. 含有简单表示 n 次的表示: I.3.1. 解析表示的逆表示: III.3.11. 点分布与函数的卷积: III.3.4. 卷积乘积: III.3.1, III.3.18. 余单位: II.1.1. 李群的 C'-连通子集: III.6.2. 判据, 卡当: I.5.4.

代数的导子: I.1.1. 李代数的内导子: I.1.2. 相对于形式群的左不变导子: I.1, 习题 24. 自由群代数中的偏导子: II.5, 习题 2. （代数）通过扩张标量由代数导出: I.1.1. （表示）通过扩张标量由表示导出: I.3.8. 导出理想: I.1.5. 导出列: I.1.5.

映入李群的映射的左微分: III.3.17, III.3.18. 表示的维数: I.3.1. 对偶表示: I.3.3.

元素, Casimir: I.3.7. 幂零: I.6.3. 李代数包络代数所关联的分次代数中的 n 次元: I.2.6. 李代数包络代数中滤过 \(\leqslant n\) 的元: I.2.6. 半单: I.6.3. 李代数中的可交换元素: I.1.3. 消去定理: II.2.9. 恩格尔定理: I.4.2. 李代数的包络代数: I.2.1. 包络双代数: II.1.4. 李代数的等价扩张: I.1.7. 穷尽滤过: II.4.1. 指数: II.6.1, III.4.3, III.6.4. 扩张, 中心: I.1.7. 非本质: I.1.7. g-模的扩张, g 的表示的扩张: I.7.2. 李代数被李代数扩张: I.1.7. 平凡: I.1.7. 扩张, 等价: I.1.7.

忠实表示: I.3.1. 点分布的域: III.3.5, III.3.18. 滤过双代数: II.1.3. 群上的实滤过: II.4.1. 第一类的典范坐标图: III.4.3. 第一类的典范坐标系: III.4.3. 与李子代数相关的左叶层: III.4.1. 形式, 双线性, 与 g-模（g 的一个表示）相关: I.3.6. 双线性, 完全不变: I.3.6. 双线性, 不变: I.3.6. Killing: I.3.6. 公式, Jacobson: I.1, 习题 19. 自由李代数: II.2.2. 李 p-代数: II.3, 习题 4. 原群: II.2.1. 函数, 与滤过相关的阶: II.4.2.

与滤过群相关的分次群: II.4.3. 相关李代数: II.4.4. L(I) 的总分次: II.2.6. 群的交换形式: I.1, 习题 24. 形式群: I.1, 习题 24. 形式线性群: I.1, 习题 25. 形式正交群: I.1, 习题 26. 形式辛群: I.1, 习题 26. 李群芽: III.1.10. 由李代数定义的李群芽: III.4.2.

Hall 基: II.2.11. 公式: II.5, 习题 9. 集: II.2.10. Hausdorff 函数: II.7.2, II.8.3. 群: II.6.2. 逆元公式: II.6, 习题 4. 级数: II.6.4. 李代数的全自同构群: I.1, 习题 16. 同态, 代数: I.1.1. 从 g 的李子代数的包络代数到 g 的包络代数的典范同态: I.2.3. 从 g 的对称代数到与 g 的包络代数相关的分次代数上的典范满同态: I.2.6. 形式同态: I.1, 习题 24.

理想, 特征: I.1.4. 导出: I.1.5. 代数的左（右、双边）理想: I.1.1. 李代数的理想: I.1.4. 恒等式, Jacobi: I.1.2. 诱导的李群结构: III.4.5. 非本质扩张: I.1.7. 无穷小自同构: III.10.1. 内导子: I.1.2. 群的整滤过: II.4.1. 李群的子群: III.6.2. 不变双线性型: I.3.6. 左不变点分布域: III.3.6. g-模（一个表示）的不变: I.3.5. 不变截面: III.3.13.

逆像李群结构: III.1.9. 不可约表示: I.3.1. 同构表示: I.3.1. g 的对称代数到其包络代数底层向量空间上的典范同构: I.2.7. g-模的同型分量: I.3.1.

Jacobi 恒等式: I.1.2. Jordan 定理: III.4, 习题 11.

扩张的核: I.1.7. Killing 型: I.3.6.

g-模（表示）的最大幂零性理想: I.4.3. 李代数的最大幂零理想: I.4.4. 运算律片段: III.1.11. 形式群的运算律: I.1, 习题 24. 无穷小运算律: III.3.7, III.3.18. 自由原群元素的长度: II.2.1. Levi 子代数: I.6.8. Levi-Malcev 定理: I.6.8. 李代数: I.1.2. 交换李代数: I.1.3. 特征幂零李代数: I.4, 习题 19. 幂零李代数: I.4.1. 形式群的李代数: I.1, 习题 24. 李群的李代数: III.3.7. 李群芽的李代数: III.3.18. 模的端同态代数: I.1.2. 约化李代数: I.6.4. 半单李代数: I.6.1. 单李代数: I.6.2. 可解李代数: I.5.1. 李群: III.1.1. Lie 定理: I.5.3. 局部同构的群芽: III.1.10. 对数: II.6.1, III.7.6. 下中心列: I.1.5.

Magnus 代数: II.5.1. Magnus 群: II.5.2. 从李代数到其包络代数的典范映射: I.2.1. Maurer-Cartan 公式: III.3.14, III.3.18.

Möbius 函数: II, 附录. 反演公式: II, 附录. 态射, 李群: III.1.2. 李群芽: III.1.10. L(I) 的多重分次: II.2.6.

幂零李代数: I.4.1, II.2.7. 李代数的根基: I.5.3. 可赋范代数: II.7. 正规化子: I.1.4, III.9.4. 赋范李代数: II.7, II.8.2.

相反代数: I.1.1. 滤过群中元素的阶: II.4.2.

p-进李群: III.1.1, III.8.1. 李 p-代数: I.1, 习题 20. p-幂单李代数: I.4, 习题 23. 交换李 p-代数的 p-核: I.1, 习题 23. p-导子: I.2, 习题 7. 幂零群的 P-包络: II.4, 习题 15. p-同态: I.1, 习题 20. p-理想: I.1, 习题 22. P-整数: II.4, 习题 14. p-映射: I.1, 习题 20. p-多项式: I.7, 习题 5. 幂零群子群的 P-饱和: II.4, 习题 14. P-挠群, P-无挠群: II.4, 习题 14. 可交换元素: I.1.3, III.9.3. Poincaré-Birkhoff-Witt 定理: I.2.7. 李多项式: II.2.4. 多项式映射: II.2.4. 李代数的呈示: II.2.3. 双代数的本原元: II.1.2. 乘积代数: I.1.1. 李群的乘积: III.1.4. 李代数的半直积: I.1.8. 表示的张量积: I.3.2, III, 附录. 种类 N 的纯 g-模: I.3.1. 纯表示: I.3.1.

李拟子群: III.1.3.

商代数: I.1.1. 李群的商: III.1.6. 商表示: I.3.1.

李代数的幂零根基: I.5.3. 李代数的根基: I.5.2. 李群的根基: III.9.7.

实李群: III.1.1, III.8.1, III.8.2.

约化李代数: I.6.4. 李代数中的（李子代数）: I.6.6.

关系子: II.2.3.

端同态的复本: I.5, 习题 14.

表示, 伴随: I.3.1. 李群的解析线性表示: III.1.2. 完全可约: I.3.1. 含有简单表示 n 次: I.3.1. 对偶: I.3.3. 忠实: I.3.1. 不可约: I.3.1. 通过扩张标量环得到的: I.3.8. 李代数的: I.3.1. 种类 σ 的纯表示: I.3.1. 商: I.3.1. 半单: I.3.1. 简单: I.3.1.

表示, 同构: I.3.1. 相似: I.3.1.

标量限制, 由李群导出的李群: III.1.1.

可解李代数的根: III.9, 习题 17.

第二类的典范坐标图: III.4.3. 第二类的典范坐标系: III.4.3.

向量丛的截面: III.1.8.

李代数的半直积: I.1.8. 李群的半直积: III.1.4.

半单李代数: I.6.1. 半单李群: III.9.8. 半单表示: I.3.1.

分离滤过: II.4.1.

列, 连接两个子代数的合成列: I.1, 习题 14. 导出列: I.1.5. 李形式幂级数: II.6.3. 李代数的下中心列: I.1.5, II.2.7. 李代数的上中心列: I.1.6.

相似表示: I.3.1. 半单李代数的简单分量: I.6.2. 简单李代数: I.6.2. 简单表示: I.3.1. 可解李代数: I.5.1. 取值于 g-模的上同调空间: I.3, 习题 12. 李齐性空间: III.1.6. 特殊自同构: I.6.8. 纯 g-模的同型分量的种类: I.3.1. 标准群: III.7.3. 子代数: I.1.1. 李子代数: I.6.8. 李代数中的约化子代数: I.6.6. 次不变子代数: I.1, 习题 14. 李子群芽: III.1.10. 李子群: III.1.3. 子表示: I.3.1. 表示的直和: I.3.1. 模的对称代数: I.2.5.

t 次幂映射: III.4.3. 复合的切运算律: III.2.1. 李子群: III.4.5. 定理, 阿多: I.7.3. 恩格尔: I.4.2. Levi-Malcev: I.6.8. Lie: I.5.3. Poincaré-Birkhoff-Witt: I.2.7. Weyl: I.6.2. Zassenhaus: I.7.2. 传递子: I.1, 习题 1. 群, 严格上三角: II.4.6. 平凡扩张: I.1.7. g-模: I.3.1. 向量 G-丛: III.1.8. T(G) 的右（分别左）平凡化: III.2.1, III.2.2. 双边理想, I.1.1. 类型 (N) 的实李群: III.9, 习题 29.

\(u\) -余代数的 u-本原元: II.1.1.\\
幂单端同态: III.9.5.\\
连通李群的万有覆盖: III.1.9.

群的上中心列: II.4, 习题 18. 李代数的中心列: I.1.6.

Weyl 定理: I.6.2.

Zassenhaus 定理: I.7.2.

